\chapter{Desarrollo}
\label{ch:desarrollo}

Como se mencionó en la introducción, el desarrollo de este ejercicio práctico consiste en la migración del
código empleado en el Ejercicio N° 2 hacia el estándar CMSIS, y la incorporación de una gestión de tiempos
precisa utilizando el temporizador de sistema (\textit{SysTick}). El comportamiento general del firmware —lectura
del ADC, manejo de interrupciones externas y control de una barra de LEDs— se mantiene inalterado, pero la forma
en la que se interactúa con el hardware subyacente se modifica por completo.

A continuación, se detallan los aspectos clave de la refactorización del código, el uso de las estructuras
provistas por el estándar, y la configuración del temporizador de sistema.

\section{Migración a la Biblioteca CMSIS}
\label{sec.ej3:cmsis}

El estándar CMSIS proporciona una capa de abstracción de hardware (HAL) independiente del fabricante para los
núcleos de procesador Cortex-M. Al incluir la cabecera \lstinline{stm32f103xb.h}, se tiene acceso a una
definición exhaustiva y estandarizada del mapa de memoria del microcontrolador.

\subsection{Uso de Estructuras y Punteros Tipados}

En la implementación anterior (\textit{Bare Metal puro}), el acceso a los registros de periféricos se
realizaba mediante la definición manual de macros y el casteo directo de direcciones de memoria a
punteros enteros volátiles (\lstinline{*(volatile uint32_t *) 0x40010800}). Este enfoque es propenso a
errores de direccionamiento y dificulta la comprensión del código.

Con la adopción de CMSIS, los registros de los periféricos se agrupan lógicamente en estructuras
de lenguaje C (\lstinline{struct}). Por ejemplo, el periférico GPIO se define a través del tipo de
dato \lstinline{GPIO_TypeDef}. Esta estructura contiene miembros que representan exactamente los registros
disponibles (\lstinline{CRL}, \lstinline{CRH}, \lstinline{ODR}, \lstinline{BSRR}, etc.), ubicados
secuencialmente en la memoria según el desplazamiento especificado en la hoja de datos.

El acceso a un puerto específico se realiza a través de un puntero predefinido que apunta a la dirección
base del periférico en el mapa de memoria. De esta forma, configurar el puerto A es tan legible como:
\begin{lstlisting}[language=C]
GPIOA->CRL |= GPIO_CRL_MODE1_0 | GPIO_CRL_MODE1_1;
\end{lstlisting}

Este paradigma garantiza que el compilador verifique los tipos de datos y los desplazamientos, brindando
mayor seguridad en tiempo de compilación.

\subsection{Uso de Máscaras Predefinidas}

Otra ventaja significativa de la migración es el uso de máscaras de bits predefinidas para la manipulación
de registros. CMSIS define constantes simbólicas para todos los campos de bits dentro de los registros.
Estas constantes, como \lstinline{RCC_APB2ENR_IOPCEN} o \lstinline{ADC_CR2_ADON}, describen semánticamente
la función del bit o campo que se está modificando.

El código resultante reemplaza los "números mágicos" empleados en la versión anterior por operaciones
a nivel de bits lógicas y autodescriptivas. El código completo portado a CMSIS se observa en el listado~\ref{list.ej3:main}.

\lstinputlisting[language={C}, caption={Código fuente principal refactorizado con CMSIS e implementación del \textit{SysTick} (\texttt{main.c}).}, label={list.ej3:main}]{sources/src/main.c}

\section{Gestión de Tiempos mediante SysTick}
\label{sec.ej3:systick}

En versiones anteriores, los retardos (\textit{delays}) se implementaban mediante la ejecución de un bucle
\lstinline{for} vacío. Este método de espera activa presenta múltiples desventajas: el tiempo
resultante depende íntimamente de la velocidad de reloj del sistema, de las optimizaciones aplicadas por
el compilador, y no ofrece un período de tiempo real determinista.

Para solucionar este inconveniente, se ha configurado el \textit{SysTick}, un temporizador de 24 bits
integrado directamente en el núcleo ARM Cortex-M3.

\subsection{Configuración del Temporizador}

El estándar CMSIS provee la función \lstinline{SysTick_Config(uint32_t ticks)}, la cual inicializa el
temporizador y habilita su interrupción de desbordamiento. Para lograr una base de tiempo de \qty{1}{\ms},
se debe cargar el registro de recarga con la cantidad equivalente de ciclos de reloj. Esto se logra
dividiendo la frecuencia del núcleo (\lstinline{SystemCoreClock}) por $1000$:

\begin{lstlisting}[language=C]
SystemCoreClockUpdate();
SysTick_Config(SystemCoreClock / 1000);
\end{lstlisting}

\subsection{Manejador de Interrupción e Implementación del Retardo}

Una vez configurado, el \textit{SysTick} genera una excepción asíncrona exactamente cada milisegundo. En
la tabla de vectores se ha implementado la rutina de servicio de interrupción (ISR) \lstinline{SysTick_Handler},
la cual incrementa una variable global volátil \lstinline{ticks}.

Esta variable se emplea para construir una función de retardo bloqueante determinista, \lstinline{delay_ms},
que calcula la diferencia de tiempos (en milisegundos reales) y cede la ejecución hasta que haya
transcurrido el período deseado. Esto permite reemplazar los retardos empíricos del bucle principal y garantizar
una ejecución periódica del parpadeo del LED interno.
